\section{Descriptive Statistics} \label{sec_descriptives}

\subsection{Distribution of project outcomes}

The resulting dataset contains \gn{NumberOfCases} agenda items with motions voted on by the CPC. Table \ref{tab_result_unanimity} shows the distribution of the project outcomes by the unanimity of the vote. Two important facts emerge about the CPC process. First, a minuscule number of projects are denied outright (\gn{CasesDeniedPct}\% of all cases). Rather than being denied, a more common occurrence is that the decision is postponed to a later date (\gn{CasesContinuedPct}\% of cases) or the requested actions are approved with major or minor conditions or modifications attached (\gn{CasesApprovedMajorPct}\% and \gn{CasesApprovedMinorPct}\% repsectively).\footnote{Examples of minor conditions and modifications include minor adjustments to the approved Floor Area Ratio, minor changes in number of approved parking spaces, and adding required security camera locations. Examples of major conditions and modifications include major changes to the allowed number of dwelling units, increases to the required number of affordable set-asides, and changes to the requested zoning designation.} \gn{CasesApprovedPct}\% of the cases were approved without any further conditions or modifications. 

Second, most decisions were unanimous (\gn{CasesUnanimousPct}\% of all cases). Because of these two facts, we do not view disagreement \emph{within} the board as a significant source of friction in moving development projects through the pipeline. Moreover, few projects are denied outright, so the impact of CPC hearings on final outcomes must come through either i) a slowing down of the process due to having to wait for the decision, which itself may be delayed multiple times; or ii) changes to the project plan that potentially add cost or time to the development, or which could  dissuade the developer from even moving forward with the project. Our primary analysis will therefore focus on the bureaucratic factors which lead to either delayed decision making or attaching conditions to the project proposal.

\subsection{Distribution of public comments}




